\documentclass[11pt,a4paper]{article}
\usepackage[utf8]{inputenc}
\usepackage[T1]{fontenc}
\usepackage[spanish]{babel}
\usepackage{hyperref}
\usepackage{geometry}
\geometry{margin=2.5cm}
\usepackage{parskip}

\title{342 se triangula con lo pendejo, no con el estado del arte}
\author{Alejandro Valenzuela Rivera\\\small Universidad Autónoma de Nuevo León}
\date{\today}
\begin{document}
\maketitle

\begin{quote}
\textbf{Sesgo declarado.} Este trabajo se sostiene desde el sesgo declarado de su autor, y no finge neutralidad. Lo que afirma se puede tumbar, y ese es el punto: el falsificador está escrito, no escondido.
\end{quote}

\section*{Resumen}
No se triangula solo con el estado del arte: se triangula con la caricatura que nadie tomó en serio y con el libro más rancio, porque entran por la misma puerta.

\section{Por qué se sostiene}
Un método que solo funciona con lo complejo no demostró nada: la respuesta es 'eso es obvio'.

\section{De dónde viene}
La autoridad: 'desde Peppa Pig hasta los libros mas rancios del mundo'.

\section{El caso}
Peppa Pig: la caricatura que nadie se tomó en serio, y entra por la misma puerta.

\section{Método}
\textquestiondown{}Se probó con lo pendejo, o solo con lo que ya estaba aceptado?

\section{Límites, y qué lo tumba}
Un caso donde el método solo funcione con lo complejo.

\section{Lo que queda abierto}
\textquestiondown{}El método aguanta la caricatura?

\section*{Cómo se leyó esta tesis}
Esto se escribió en el repositorio man-made-empathy, del que este candado es una pieza: candados/06 del caso/342 se triangula con lo pendejo, no con el estado del arte.md. La máquina que lo compuso lo hizo desde este mismo archivo, y no agregó nada que no estuviera escrito. Si algo de esto es falso, se dice en el mismo sitio, y se corrige con fecha.

\end{document}
